\documentclass[12pt]{article}
\usepackage[utf8]{inputenc}
\usepackage[T1]{fontenc}
\usepackage[margin=1in]{geometry}
\usepackage{setspace}
\usepackage{xcolor}
\usepackage[colorlinks=true, urlcolor=blue!60!black, linkcolor=blue!60!black]{hyperref}
\setlength{\parindent}{0em}
\setlength{\parskip}{0.45em}
\renewcommand{\baselinestretch}{1.0}
\begin{document}
\begin{center}
{\Large \textit{Central Hudson Gas \& Electric Corp. v. Public Service Commission}, 447 U.S. 557 (1980)\par}
\vspace{0.3em}
{\normalsize Case Brief --- One-Page Edition\par}
\end{center}

\textbf{Memory Jogger.} The four-part commercial speech test; New York's utility ad ban fails (\href{https://supreme.justia.com/cases/federal/us/447/557/}{Justia}).

\textbf{Facts.} After the 1973 oil embargo, New York's Public Service Commission banned electric utilities' advertising ``promot[ing] the use of electricity,'' later citing conservation and rate fairness. Central Hudson, a regulated monopoly, challenged the ban.

\textbf{Procedural History.} The New York courts upheld the ban; the Supreme Court reversed.

\textbf{Votes.} 8--1: Powell for the Court; Brennan, Blackmun, and Stevens concurred in the judgment; Rehnquist dissented.

\textbf{Holding.} Restrictions on truthful commercial speech about lawful activity must serve a substantial interest, directly advance it, and be no more extensive than necessary; the total ban was too broad.

\textbf{Reasoning.} Conservation was substantial and directly linked to advertising, but the ban reached ads for efficient devices such as heat pumps, and narrower measures were available. Blackmun called suppression of truthful information to ``dampen'' demand an attack that ``strikes at the heart of the First Amendment''; Rehnquist saw a revival of \textit{Lochner}.

\textbf{Examples.} \emph{Same-Side:} a ban on liquor price ads. \emph{Opposite-Side:} ads for illegal betting. \emph{Fence-Sitter:} state-legal marijuana billboards.

\textbf{Critique.} Two inconsistent definitions of commercial speech; fears of First Amendment Lochnerism.

\textbf{Key Quotation.} ``In commercial speech cases, then, a four-part analysis has developed. At the outset, we must determine whether the expression is protected by the First Amendment. For commercial speech to come within that provision, it at least must concern lawful activity and not be misleading. Next, we ask whether the asserted governmental interest is substantial. If both inquiries yield positive answers, we must determine whether the regulation directly advances the governmental interest asserted, and whether it is not more extensive than is necessary to serve that interest.'' --- Justice Powell, 447 U.S. at 566.

\end{document}
